\chapter[Finite subextensions]{From Finite Subextensions to an Algebraic Galois Extension}
\label{ch:finite-subextensions}

Fix a field extension \(L/K\). We identify \(K\) with its image in
\(L\). An intermediate field is a subfield of this particular \(L\)
containing \(K\), and its scalar structures are the inclusions. If an
abstract extension of \(K\) is given with a \(K\)-embedding into \(L\),
we use its image as an intermediate field and transport its operations
through that embedding. For \(a\) algebraic over a field \(F\), write
\(\mu_{a,F}\) for the monic minimal polynomial of \(a\) over \(F\).

\section{Finite results and passage to an intermediate field}

\begin{proposition}[Finite ingredients]\label{res:F01}
Let \(a\) be algebraic over \(K\), and let \(T\) be an extension of
\(K\). The \(K\)-embeddings \(K(a)\to T\) correspond to the roots of
\(\mu_{a,K}\) in \(T\), by evaluation at \(a\). The field generated by
all the roots of a separable polynomial in a splitting extension is
finite Galois over \(K\).

If \(E/K\) is finite Galois, then
\[
 M\longmapsto\operatorname{Gal}(E/M),\qquad
 H\longmapsto E^H
\]
give inverse inclusion-reversing maps between the intermediate fields
of \(E/K\) and the subgroups of \(\operatorname{Gal}(E/K)\).
In particular \(E^{\operatorname{Gal}(E/K)}=K\).
For an intermediate field \(M\) Galois over \(K\), restriction
\(\operatorname{Gal}(E/K)\to\operatorname{Gal}(M/K)\) is surjective
with kernel \(\operatorname{Gal}(E/M)\).
\end{proposition}

These are the finite results used below. For the first statement the
presentation \(K(a)=K[X]/(\mu_{a,K})\) explains the root parametrization
and shows that the embedding associated with a given root is unique.
The splitting field facts, the finite correspondence, and the finite
restriction sequence are established in \cite[Chapters 12 and 14]{d07}.
See also \cite[Chapters 2 and 3]{milne} and
\cite[Tag 09DW]{stacks}. When counting embeddings of a finite
separable extension \(M/K\) into an extension containing the conjugates
of all elements of \(M\), we use a primitive element \(a\) of \(M/K\).
Its minimal polynomial then has \([M:K]\) distinct roots in that
extension, and the first statement gives the same number of embeddings.

We call \(L/K\) Galois if it is algebraic, normal, and separable.
Normality means that every irreducible polynomial in \(K[X]\) having
a root in \(L\) splits in \(L[X]\). In an algebraic extension separability
can be checked on the minimal polynomial of every element. We write
\[
 G=\operatorname{Gal}(L/K),\qquad G_M=\operatorname{Gal}(L/M)
\]
for the group of \(K\)-automorphisms and the subgroup fixing an
intermediate field \(M\) pointwise. If \(M/K\) is Galois, put
\(Q_M=\operatorname{Gal}(M/K)\). For a subgroup \(H\) of \(G\), \(L^H\)
denotes the field \(\{x\in L:h(x)=x\ \text{for every }h\in H\}\).

\begin{proposition}[The upper part of a Galois tower]\label{res:F02}
If \(L/K\) is Galois and \(M\) is an intermediate field, then \(M/K\)
is algebraic and separable, and \(L/M\) is Galois.
\end{proposition}
\begin{proof}
For \(a\in M\), the polynomial \(\mu_{a,K}\) is the same whether \(a\)
is viewed in \(M\) or in \(L\). This proves algebraicity and separability
for \(M/K\). For every \(b\in L\), the polynomial \(\mu_{b,K}\),
regarded as a polynomial over \(M\), annihilates \(b\). Thus \(b\)
is algebraic over \(M\), and \(\mu_{b,M}\) divides \(\mu_{b,K}\) in
\(M[X]\). Over \(L\) the latter is a product of distinct linear factors.
Each of its monic divisors over \(L\) therefore has the same splitting
and distinctness properties. This proves normality and separability
of \(L/M\). Compare \cite[Proposition 7.3]{milne} and
\cite[Tag 09I2]{stacks}.
\end{proof}

\section{Finite Galois subextensions}

\begin{proposition}[A finite Galois subextension containing a finite set]
\label{res:F03}
If \(L/K\) is Galois, every finite subset \(S\subseteq L\) lies in an
intermediate field \(E\) which is finite Galois over \(K\).
\end{proposition}
\begin{proof}
The empty subset lies in \(K\). Otherwise, take the product of the
distinct polynomials \(\mu_{a,K}\) for \(a\in S\). Each factor is
separable and splits in \(L\), and two distinct monic irreducible
factors have disjoint root sets. The product is therefore separable.
Inside \(L\) take the field \(E\) generated by all of its roots.
This field contains \(S\), and it is precisely the splitting field
of the product inside \(L\). Proposition~\ref{res:F01} makes it finite
Galois over \(K\). The choice of distinct factors also applies if
different elements of \(S\) are conjugate. See \cite[Chapter 7, p.~95]{milne}
and \cite[Tag 0BMG]{stacks}.
\end{proof}

\begin{theorem}[The directed union of the finite Galois subextensions]
\label{res:F04}
In a Galois extension \(L/K\), the finite Galois intermediate fields,
ordered by inclusion, form a nonempty directed set. The compositum of
two of them inside \(L\) is again one of them, and their union is \(L\).
Conversely, if an extension of \(K\) is the union of a nonempty directed
family of finite Galois intermediate fields over \(K\), the extension
is Galois.
\end{theorem}
\begin{proof}
The ground field \(K\) is a finite Galois intermediate field.
Let \(E_1,E_2\) be two such fields. In each \(E_i\) choose a finite
set of field generators over \(K\), for example a \(K\)-basis.
The minimal polynomials of these generators over \(K\) split in
\(E_i\). The field generated by their roots in \(L\) is contained
in \(E_i\) and contains all its generators, so it equals \(E_i\).
Thus \(E_1E_2\) is the field generated by the roots of the polynomials
from both lists. It is the splitting field of their product with
duplicate monic irreducible factors omitted. That product is
separable, and Proposition~\ref{res:F01} makes \(E_1E_2/K\) finite
Galois. This provides a common upper bound for \(E_1\) and \(E_2\).
Proposition~\ref{res:F03} applied to singletons proves the union
assertion.

For the converse, directedness ensures that two elements of the
union lie in a common field, so the union is a field. Let \(a\)
lie in the union. One field in the family contains \(a\), and
its being finite over \(K\) proves that
\(a\) is algebraic over \(K\). The minimal polynomial \(\mu_{a,K}\)
is separable and splits in that same field. This holds for every
\(a\) and proves both separability and normality of the union over
\(K\). See \cite[Proposition 3.21, Chapter 7, and Example 7.26]{milne},
\cite[Theorem 2.11.1]{rz}, and \cite[Tag 0BU2]{stacks}.
\end{proof}

\section{Extension of embeddings and restriction}

\begin{theorem}[Extension of an intermediate-field embedding]
\label{res:F05}
Suppose \(L/K\) is Galois. If \(M\) is an intermediate field and
\(\iota:M\to L\) is a \(K\)-embedding, there is a \(K\)-automorphism
\(\sigma:L\to L\) with \(\sigma|_M=\iota\).
\end{theorem}
\begin{proof}
Choose an algebraic closure \(\Omega\) of \(L\), using the given
inclusion \(L\subseteq\Omega\). Consider the pairs \((N,\tau)\) with
\(M\subseteq N\subseteq L\) and \(\tau:N\to\Omega\) a \(K\)-embedding
extending \(\iota\) as a map into \(\Omega\). Order them by extension
of both domain and map. They include \((M,\iota)\). Along a chain
the union of the domains is a field and the embeddings agree on
overlaps. Their union is an embedding, since any pair of arguments
lies in a common member of the chain. Hence every nonempty chain
has an upper bound, and Zorn's lemma gives a maximal pair.

If its domain \(N\) is proper in \(L\), choose \(x\in L\setminus N\).
By algebraicity over \(K\), \(x\) is algebraic over \(N\).
Transport \(\mu_{x,N}\) along \(\tau\); the result is irreducible
over \(\tau(N)\) and has a root in \(\Omega\).
The quotient presentation of a simple extension in
Proposition~\ref{res:F01} extends \(\tau\) to \(N(x)\) by sending
\(x\) to that root. This contradicts maximality. We obtain a
\(K\)-embedding \(\tau:L\to\Omega\).

For \(x\in L\), its image \(\tau(x)\) is a root of \(\mu_{x,K}\).
Normality of \(L/K\) puts all the roots of this polynomial in \(L\),
so we can regard \(\tau\) as a \(K\)-endomorphism of \(L\).
It is surjective: the finite set of roots of \(\mu_{x,K}\) is
stable under \(\tau\), and the restriction to this set is injective,
hence a permutation. In particular \(x\) has a preimage in \(L\).
This is true for every \(x\), so \(\tau\) is the required
automorphism. See \cite[Theorem 6.6 and Proposition 7.4]{milne},
\cite[Chapter 12]{d07}, and \cite[Tag 0BME]{stacks}.
\end{proof}

\begin{corollary}[Stable intermediate fields]\label{res:F06}
In a Galois extension \(L/K\), an intermediate field \(M\) is Galois
over \(K\) if and only if \(\sigma(M)=M\) for every \(\sigma\in G\).
\end{corollary}
\begin{proof}
If \(M/K\) is Galois and \(a\in M\), then \(\sigma(a)\) is a root of
\(\mu_{a,K}\), whose roots all lie in \(M\). Thus \(\sigma(M)\subseteq M\).
Apply this also to \(\sigma^{-1}\) to get equality.
Conversely, suppose \(M\) is stable under \(G\). For any \(a\in M\)
every root \(b\) of \(\mu_{a,K}\) in an algebraic closure of \(L\)
already lies in \(L\). By Proposition~\ref{res:F01}, \(a\mapsto b\)
defines a \(K\)-embedding \(K(a)\to L\).
Theorem~\ref{res:F05} extends it to a member of \(G\).
Stability of \(M\) gives \(b\in M\), so \(\mu_{a,K}\) splits in \(M\).
This proves normality over \(K\). Algebraicity and separability
follow from Proposition~\ref{res:F02}. See
\cite[Corollary 7.5 and Theorem 7.13(d)]{milne}.
\end{proof}

\begin{proposition}[Surjective restriction and composition]
\label{res:F07}
Suppose \(L/K\) is Galois. If \(M/K\) is a Galois intermediate
extension, restriction gives a surjective homomorphism
\[
 r_M:G\longrightarrow Q_M,\qquad r_M(\sigma)=\sigma|_M,\qquad
 \ker r_M=G_M.
\]
For any two Galois intermediate fields \(E\subseteq E'\),
restriction also gives a surjective homomorphism
\(r_{E',E}:Q_{E'}\to Q_E\).
These maps satisfy \(r_{E,E}=\operatorname{id}_{Q_E}\),
\(r_E=r_{E',E}\circ r_{E'}\), and
\[
 r_{E'',E}=r_{E',E}\circ r_{E'',E'}
 \quad\text{for } E\subseteq E'\subseteq E''
\]
among Galois intermediate fields.
\end{proposition}
\begin{proof}
Corollary~\ref{res:F06} ensures that \(\sigma|_M\) maps \(M\) onto
itself, so restriction takes values in \(Q_M\).
The restriction of a composite agrees with the composite of the
restrictions, since every restricted map has domain and image \(M\).
An element of \(Q_M\), composed with the inclusion of \(M\) into \(L\),
is a \(K\)-embedding \(M\to L\), and Theorem~\ref{res:F05} extends
it to a member of \(G\). This proves surjectivity. The kernel is
precisely the set of maps acting identically on \(M\).
For \(r_{E',E}\) apply the same argument with \(E'\) as top field.
Each identity of restriction maps follows by evaluating both
sides on an element of the smallest field involved. In particular
these identities make the finite-stage restriction maps an inverse
system. The continuity assertions for restriction will use the
topology constructed in the next chapter. See
\cite[Proposition 7.7 and Example 7.26]{milne} and
\cite[Tag 0BMK]{stacks}.
\end{proof}
